% ===============================================
% (13) 국문초록 — placed at the BACK, after the appendix
% Maximum 2 pages. The guideline requires this page to be
% written in Korean, so kotex is used for hangul rendering.
% ===============================================
\clearpage
\thispagestyle{plain}
\addstarchaptertoc{국문초록}

\begin{center}
{\fontsize{16pt}{22pt}\selectfont\bfseries 국문초록\par}

\vspace{8mm}

{\fontsize{14pt}{20pt}\selectfont\bfseries \ThesisTitleKR\par}
\end{center}

\vspace{6mm}

\begin{flushright}
{\fontsize{12pt}{18pt}\selectfont
\AuthorKR\par
\DepartmentKR\par
\GraduateSchoolKR\par}
\end{flushright}

\vspace{8mm}

% Body — 11 pt, line spacing 1.5 (or 2.0), 2-character indent
\setstretch{1.5}

여기에 국문초록 본문을 입력하세요. 본문은 11\,pt, 줄간격 1.5 또는 2.0,
문단 시작마다 2칸 들여쓰기를 적용합니다. 본 템플릿은 가천대학교 일반대학원의
영문 학위논문 작성지침(Thesis/Dissertation Guidelines)을 근거로 작성되었으며,
영문 지침에 따라 국문초록은 본문 뒤에 배치됩니다.

국문초록은 연구의 배경과 필요성, 목적, 방법, 주요 결과, 결론을 2쪽 이내로
간결하게 기술합니다. 본문 마지막에는 핵심어를 7$\sim$8단어 이내로 기입합니다.

\vfill

\noindent\textbf{핵심어}: 키워드1, 키워드2, 키워드3, 키워드4, 키워드5, 키워드6, 키워드7

\clearpage
